\chapter{\texorpdfstring{\uppercase{AI-Assisted Prompts}}{AI-Assisted Prompts}}
\label{app:ai-prompts}

\newsavebox{\aipromptbox}
\newenvironment{aiprompt}
{\par\medskip\noindent\begin{center}\begin{lrbox}{\aipromptbox}\begin{minipage}{0.96\textwidth}\sffamily\normalsize\setlength{\parindent}{0pt}\setlength{\parskip}{0.45em}}
{\end{minipage}\end{lrbox}\setlength{\fboxsep}{9pt}\fbox{\usebox{\aipromptbox}}\end{center}\medskip}

This appendix records representative prompts corresponding to the limited uses of
AI tools declared in the report. The prompts are provided for transparency and
reproducibility. They define the intended scope of assistance and require the
author to verify all scholarly claims, sources, and editorial changes before
incorporation into the report.

\section{\texorpdfstring{\uppercase{Literature-Survey Support}}{Literature-Survey Support}}

\noindent\textbf{Prompt B.1: Identification of relevant scholarly sources}

\begin{aiprompt}
Identify peer-reviewed papers, authoritative conference publications, or
official technical documentation relevant to the following literature-survey
topics: \textit{context paging for agent memory and optimized planning}. For each suggested source,
provide the complete bibliographic details, publication venue and year, DOI or
stable URL where available, and a concise explanation of its relevance. Do not
invent sources, bibliographic information, results, or quotations. Distinguish
clearly between peer-reviewed research, preprints, and technical documentation.
Suggest only sources that can be independently verified by the author. Do not
write unsupported claims for direct inclusion in the report; instead, identify
which claims each verified source can reasonably support and note any important
limitations or differences from the present project.
\end{aiprompt}

\section{\texorpdfstring{\uppercase{Language and Editorial Support}}{Language and Editorial Support}}

\noindent\textbf{Prompt B.2: Academic language and grammar review}

\begin{aiprompt}
Review the following report passage for grammar, spelling, punctuation,
sentence structure, clarity, consistency, and formal academic style. Preserve
the author's technical meaning, numerical results, citations, terminology,
section structure, and stated limitations. Do not add new claims, references,
experimental results, or interpretations. Identify any sentence whose meaning
is ambiguous and explain the issue before proposing a revision. Return: (1) a
revised passage with only necessary editorial changes, and (2) a concise list of
the substantive edits made.
\end{aiprompt}
\clearpage
\section{\texorpdfstring{\uppercase{Synthetic Test-Case Support}}{Synthetic Test-Case Support}}

\noindent\textbf{Prompt B.3: Development of candidate synthetic test cases}

\begin{aiprompt}
Using the following approved evaluation objective and system description,
propose a set of deterministic synthetic test cases for evaluating
\textit{context paging for agent memory, priority-aging scheduling of multi-agent systems, and optimized planning}. Each case
should specify a unique identifier, the input scenario, the relevant context
or task information, the expected system behaviour, and the measurable success
criteria. Ensure that the cases cover both normal operation and meaningful
edge conditions, such as missing context, irrelevant retrieved context,
repeated information, dependency constraints, and invalid or ambiguous plans.
Do not invent implementation capabilities that are not stated in the system
description. Do not include private, sensitive, or personally identifying
information. Avoid duplicate cases and avoid wording that reveals the expected
answer unless it is required for deterministic evaluation. Return the cases in
a structured table and clearly identify assumptions.
\end{aiprompt}

